\section{논의와 한계}\label{sec:discussion}
\subsection{현재 결과가 지지하는 범위}
이미 학습된 추정기의 코너를 수정한다는 개념과 치수·대칭 활용 자체는 선행연구에 존재한다. 본 결과가 보여주는 것은 내부 영상 특징과 알려진 치수로 제한된 이동을 학습했을 때, 세 기반에서 중심부 오차가 어떻게 줄고 어떤 실패가 남는지이다. 여러 기반의 결과는 각 기반에 별도로 학습한 보정기의 적용성으로 해석해야 한다. 초기 디코더와 검출 조건까지 완전히 같게 만든 백본 치환 실험이나 하나의 가중치를 임의 추정기에 전이한 결과는 아니다.

Base 대비 이득에는 추가 보정 학습 전체의 효과가 포함된다. N0--N3 절제는 치수 분기와 회전 대응의 역할을 더 좁히지만 치수 내용과 파라미터 효과를 완전히 분리하지 못한다. 치수 입력이 후보 점수에 영향을 준다는 것과 최종 코너가 물리 형상을 일관되게 복원한다는 것도 다르다. 더 강한 형상 일관성이나 초기 추정기 직접 치수 입력의 이득은 이 후단 보정 결과로 대체할 수 없다.

\subsection{대안 선택과 실패의 의미}
N3에 포함된 회전 대응 감독의 추가 이득은 작거나 혼합되어 있다. N2를 더 단순한 후보로 논의할 수는 있으나, N3의 결과를 N2의 결과로 바꾸어 부를 수는 없다. PoseFix 기반 보정도 여러 중앙값에서 N3보다 우수하므로, P90 하나만으로 제안 방법을 전반적으로 우월하다고 결론내리지 않는다. 동일 정보·학습 선택·실행환경의 대안 비교가 더 확보되어야 사용 조건에 따른 선택을 강화할 수 있다.

심한 가림에서 위치 오차가 악화하고 큰 오류 복구가 드물다는 결과는 방법의 작동 범위를 제한한다. 출력 박스와 점수 보존은 검출 경로에 개입하지 않았다는 뜻이며, 원래 정확한 코너나 자세를 항상 보존했다는 뜻은 아니다. 코너 개선이 자세 개선으로 전달되는지 직접 확인한 이유도 여기에 있다.

\subsection{자료와 측정의 제한}
직사각형 319장은 재사용 개발자료이고 정사각형은 한 세션·한 치수이다. 직사각형 가림 등급은 사람 완료 제출본과 연결했지만, 정확한 판정 경계·예측 노출·표본 모집 이력이 모두 독립적으로 통제된 것은 아니다. 두 128장 집합의 frame ID는 전수 일치했으나, 실제 개체 수·선정 이력 및 과거 정사각형 자료와의 전체 노출 관계는 추가 확인 대상이다. 정사각형 119장의 가림 등급과 대부분의 코너 가시성도 사람 검수가 남아 있다.

자세 참조는 수동 코너와 물체 기하의 재구성에 의존한다. 개별 코너 참조의 생성 경로와 작은 차이에 대한 참조 불확실성이 충분히 정량화되지 않았으므로 독립 물리 6D 실측 정확도를 주장하지 않는다. 세션 재표집 구간은 관측 결과의 변동을 설명하지만 독립 세션 검증이나 개발 중 선택 편향 보정을 대신하지 않는다. 기록 영상 사례와 정사각형의 독립 자세 정확도는 실제 참조가 검증된 뒤 별도로 완성해야 한다.
